\section{SAT-basierter Solver} \label{sec:sat}

Dieses Kapitel beschreibt einen Sudoku-Solver, der das Rätsel als aussagenlogische Formel formuliert und mit einem \acs{SAT}-Solver löst.
Das Lösen durchläuft drei Schritte: Das Gitter wird in Klauseln kodiert, ein \acs{SAT}-Solver sucht ein Modell, und dieses Modell wird zurück in ein ausgefülltes Gitter übersetzt.
Abschnitt~\ref{sec:sat-encoding} führt \acs{SAT} ein und beschreibt die Kodierung, Abschnitt~\ref{sec:algorithms} die Suchalgorithmen \acl{DPLL} (\acs{DPLL})\acused{DPLL} und \acl{CDCL} (\acs{CDCL})\acused{CDCL} und Abschnitt~\ref{sec:implementation} deren Implementierung.

\subsection{SAT und die Sudoku-Kodierung} \label{sec:sat-encoding}

Die Solver in diesem Kapitel arbeiten auf aussagenlogischen Formeln in \emph{konjunktiver Normalform} (\acs{KNF})\acused{KNF}.
Damit ein \acs{SAT}-Solver ein Sudoku lösen kann, müssen dessen Regeln daher als \acs{KNF}-Formel vorliegen.
Abschnitt~\ref{sec:sat-cnf} führt die nötigen Grundbegriffe ein, Abschnitt~\ref{sec:encoding} übersetzt darauf aufbauend die Sudoku-Regeln in Klauseln.

\subsubsection{SAT und KNF} \label{sec:sat-cnf}

Ein \emph{Literal} ist eine boolesche Variable $x$ oder ihre Negation $\neg x$.
Eine \emph{Klausel} ist eine Disjunktion von Literalen. Sie ist erfüllt, sobald eines ihrer Literale wahr ist.
Eine \acs{KNF}-Formel ist eine Konjunktion von Klauseln \vgl[204]{davis1960}, zum Beispiel
\[
    F = (x_1 \lor \neg x_2) \land (x_2 \lor x_3) \land (\neg x_1 \lor \neg x_3).
\]
Eine Belegung mit Wahrheitswerten, die jede Klausel erfüllt, ist ein \emph{Modell} von $F$.
Hier ist $x_1 = x_2 = \text{wahr}$ und $x_3 = \text{falsch}$ ein solches.
\acs{SAT} ist die Frage, ob eine Formel ein Modell hat.

\acs{SAT} war das erste Problem, dessen NP-Vollständigkeit bewiesen wurde \vgl[151\psq]{cook1971}. Daher ist kein Algorithmus bekannt, der jede Instanz in Polynomialzeit löst.
Die beiden möglichen Antworten unterscheiden sich darin, wie schwer sie zu überprüfen sind.
Ein Modell lässt sich in linearer Zeit verifizieren, indem jede Klausel ausgewertet wird.
Für die Antwort \enquote{unerfüllbar} ist im Allgemeinen kein kurzes Zertifikat bekannt.
Die Beschränkung der Eingabe auf \acs{KNF} ist keine Einschränkung der Allgemeinheit: Jede aussagenlogische Formel lässt sich in eine erfüllbarkeitsäquivalente \acs{KNF}-Formel linearer Größe umformen \vgl[115\psq]{tseitin1970}.
Sudoku benötigt diesen Schritt nicht, da sich seine Regeln direkt als Klauseln formulieren lassen (Abschnitt~\ref{sec:encoding}).

Die Implementierung verwendet die ganzzahlige Notation des DIMACS-Formats, des gängigen Eingabeformats von \acs{SAT}-Solvern \vgl[2\psq]{dimacs1993}.
Die Variablen werden ab 1 nummeriert, das Literal $x_i$ wird als $i$ und $\neg x_i$ als $-i$ geschrieben, und eine Klausel ist eine Liste solcher ganzen Zahlen.
Aus $F$ wird damit \texttt{[[1, -2], [2, 3], [-1, -3]]}.

\subsubsection{Kodierung} \label{sec:encoding}

Die Kodierung folgt Lynce und Ouaknine~\parencite*[3\psqq]{lynce2006}.
Für jede Zelle $(r, c)$ und jede Ziffer $d$ gibt es eine boolesche Variable $x_{r,c,d}$, die genau dann wahr ist, wenn Zelle $(r, c)$ die Ziffer $d$ enthält.
Ein $n \times n$-Gitter benötigt daher $n^3$ Variablen, ein $9 \times 9$-Gitter also 729.
Die Implementierung nummeriert sie als
\[
    \mathrm{var}(r, c, d) = r \cdot n^2 + c \cdot n + d,
    \quad \text{mit } r, c \in \{0, \dots, n-1\} \text{ und } d \in \{1, \dots, n\}.
\]
Damit werden die Variablen lückenlos und eineindeutig auf $1, \dots, n^3$ abgebildet, wie es die DIMACS-Notation verlangt.
Das Dekodieren eines Modells ist die Umkehrrechnung zur Basis $n$.

Die Sudoku-Regeln werden zu vier Gruppen von Klauseln (Tabelle~\ref{tab:clause-groups}).
Zwei Gruppen fordern, dass etwas mindestens einmal gilt (\emph{\acl{ALO}}, \acs{ALO})\acused{ALO}, jeweils als eine einzige lange Klausel.
Die beiden anderen fordern, dass etwas höchstens einmal gilt (\emph{\acl{AMO}}, \acs{AMO})\acused{AMO}.
Sie werden paarweise kodiert: eine binäre Klausel für jedes Paar von Literalen, die nicht beide wahr sein dürfen.
Zusätzlich fügt jede vorgegebene Ziffer $g$ in Zelle $(r, c)$ die Unit-Klausel $(x_{r,c,g})$ hinzu.

\begin{table}[htb]
    \centering
    \caption{Klauselgruppen der Kodierung. $U$ ist eine Einheit: eine der $n$ Zeilen, $n$ Spalten oder $n$ Blöcke.}
    \label{tab:clause-groups}
    \small
    \renewcommand{\arraystretch}{1.3}
    \begin{tabularx}{\textwidth}{@{}>{\raggedright\arraybackslash}X l l r@{}}
        \toprule
        Gruppe & Klauseln & Anzahl & $9 \times 9$ \\
        \midrule
        \textbf{Zelle definiert} (\acs{ALO}): jede Zelle enthält mindestens eine Ziffer
            & $\bigwedge_{r,c} (x_{r,c,1} \lor \dots \lor x_{r,c,n})$
            & $n^2$ & 81 \\
        \textbf{Zelle eindeutig} (\acs{AMO}): jede Zelle enthält höchstens eine Ziffer
            & $\bigwedge_{r,c} \bigwedge_{d < d^{\prime}} (\neg x_{r,c,d} \lor \neg x_{r,c,d^{\prime}})$
            & $n^2 \binom{n}{2}$ & 2.916 \\
        \textbf{Einheit eindeutig} (\acs{AMO}): jede Ziffer kommt höchstens einmal pro Einheit vor
            & $\bigwedge_{U,d} \bigwedge_{(r,c) < (r^{\prime},c^{\prime}) \in U} (\neg x_{r,c,d} \lor \neg x_{r^{\prime},c^{\prime},d})$
            & $3n \cdot n \binom{n}{2}$ & 8.748 \\
        \textbf{Einheit definiert} (\acs{ALO}): jede Ziffer kommt mindestens einmal pro Einheit vor
            & $\bigwedge_{U,d} \bigvee_{(r,c) \in U} x_{r,c,d}$
            & $3n \cdot n$ & 243 \\
        \bottomrule
    \end{tabularx}
    \par\smallskip
    {\footnotesize Quelle: eigene Darstellung in Anlehnung an Lynce und Ouaknine~\parencite*[4\psq]{lynce2006}.}
\end{table}

Die \emph{minimale Kodierung} besteht nur aus den Gruppen \enquote{Zelle definiert} und \enquote{Einheit eindeutig} und ist bereits korrekt.
Eine Einheit hat $n$ Zellen, von denen jede mindestens eine Ziffer enthält, während jede der $n$ Ziffern höchstens einmal vorkommen darf.
Also enthält jede Zelle genau eine Ziffer (Schubfachprinzip).
Die beiden anderen Gruppen sind daher redundant.
Die \emph{erweiterte Kodierung} fügt sie dennoch hinzu: Sie ändern die Lösungsmenge nicht, geben dem Solver aber mehr zu propagieren.
So kann die Unit Propagation (Abschnitt~\ref{sec:dpll}) mithilfe von \enquote{Einheit definiert} eine Ziffer setzen, die nur in eine einzige Zelle einer Einheit passt. Menschliche Löser nennen dies ein \emph{Hidden Single}.

Für ein $9 \times 9$-Gitter hat die minimale Kodierung 8.829 Klauseln und die erweiterte 11.988, was mit Lynce und Ouaknine~\parencite*[4\psq]{lynce2006} übereinstimmt.
Beide wachsen mit $O(n^4)$.
Die zusätzlichen Klauseln zahlen sich aus, sobald der Solver suchen muss: Bei drei $16 \times 16$-Rätseln mit eindeutiger Lösung benötigte der \acs{CDCL}-Solver aus Abschnitt~\ref{sec:cdcl} im Median 3.134 Konflikte mit der minimalen und 71 mit der erweiterten Kodierung.

\subsection{Algorithmen} \label{sec:algorithms}

Beide Solver dieses Kapitels gehen auf den DPLL-Algorithmus zurück.
Abschnitt~\ref{sec:dpll} beschreibt DPLL, Abschnitt~\ref{sec:cdcl} dessen Erweiterung \acs{CDCL}, die aus Konflikten lernt.
Beide Algorithmen propagieren auf dieselbe Weise und unterscheiden sich nur im Umgang mit Konflikten; Abschnitt~\ref{sec:comparison} stellt sie gegenüber.

\subsubsection{DPLL} \label{sec:dpll}

Das erste Verfahren für \acs{SAT} eliminierte Variablen durch Resolution und benötigte in der Praxis zu viel Speicher \parencites[vgl.][210]{davis1960}[394]{davis1962}.
Die Ersetzung dieses Schritts durch Verzweigen mit Backtracking führte zum DPLL-Algorithmus, benannt nach Davis, Putnam, Logemann und Loveland \vgl[394]{davis1962}.

DPLL vereinfacht die Formel zunächst so weit wie möglich.
Die wichtigste Regel ist die \emph{Unit Propagation}: Sind alle Literale einer Klausel bis auf eines falsch, muss das verbleibende Literal wahr sein.
Dessen Belegung kann weitere Klauseln zu Unit-Klauseln machen, daher wird die Regel wiederholt, bis keine solche Klausel mehr übrig ist.
Optional belegt die \emph{Pure-Literal-Regel} zusätzlich jede Variable, die nur mit einem Vorzeichen vorkommt, so, dass alle ihre Vorkommen wahr sind.
Enthält danach eine Klausel nur noch falsche Literale, führt die aktuelle Belegung zu einem Konflikt.
Ist jede Klausel erfüllt, ist die Formel erfüllbar.
Andernfalls \emph{verzweigt} DPLL: Es wählt eine freie Variable $x$, setzt sie auf wahr und beginnt erneut mit der Propagation.
Führt dies zu einem Konflikt, kehrt der Solver per Backtracking zurück und versucht stattdessen $x = \text{falsch}$.
Scheitern beide Werte, wird der Konflikt an die vorherige Verzweigung weitergereicht \parencites[vgl.][209--212]{davis1960}[394\psq]{davis1962}.

Beide Vereinfachungsregeln erhalten die Erfüllbarkeit, und die Verzweigung probiert beide Werte jeder gewählten Variable aus.
Daher terminiert DPLL stets mit der korrekten Antwort; diese lautet \enquote{unerfüllbar}, sobald beide Werte der ersten gewählten Variable gescheitert sind.
Im schlechtesten Fall probiert es jedoch weiterhin alle $2^V$ Belegungen von $V$ Variablen aus.

Auf Sudoku angewandt führt die Unit Propagation die grundlegenden Züge menschlicher Löser aus.
Eine gesetzte Ziffer wird als Kandidat aus den übrigen Zellen ihrer Einheiten entfernt (\enquote{Einheit eindeutig}).
Eine Zelle mit nur noch einem Kandidaten, ein \emph{Naked Single}, wird über \enquote{Zelle definiert} gefüllt.
Ein \emph{Hidden Single}, also eine Ziffer, die nur in eine einzige Zelle einer Einheit passt, benötigt \enquote{Einheit definiert} und wird daher nur mit der erweiterten Kodierung gefunden.
Das Verzweigen entspricht dem Raten.

\subsubsection{CDCL} \label{sec:cdcl}

Gerät DPLL in einen Konflikt, macht es nur die letzte Entscheidung rückgängig und vergisst, warum der Konflikt entstanden ist.
Dieselbe Kombination von Belegungen kann daher in anderen Teilen der Suche immer wieder denselben Konflikt verursachen.
\acs{CDCL} analysiert stattdessen jeden Konflikt.
Es fügt die Ursache als neue Klausel zur Formel hinzu und springt direkt dorthin zurück, wo diese Ursache begann \vgl[509\psqq]{marquessilva1999}.
Zusammen mit den Techniken des Chaff-Solvers \vgl[530\psqq]{moskewicz2001} bildet dies die Grundlage praktisch aller modernen \acs{SAT}-Solver.

\acs{CDCL} behält die Unit Propagation von DPLL bei, speichert aber mehr Informationen.
Jede Entscheidung eröffnet eine neue \emph{Entscheidungsebene}.
Jedes danach propagierte Literal gehört zu dieser Ebene und merkt sich seinen \emph{Grund}: die Klausel, die zur Unit-Klausel wurde und es erzwungen hat.
Entscheidungen und Gründe bilden den \emph{Implikationsgraphen}: Seine Knoten sind die belegten Literale, und jedes propagierte Literal hat eingehende Kanten von den übrigen, falschen Literalen seines Grundes.
Ein Konflikt ist eine Klausel, deren Literale alle falsch sind.

Die Konfliktanalyse durchläuft diesen Graphen rückwärts, ausgehend von der Konfliktklausel.
In jedem Schritt resolviert sie die aktuelle Klausel mit dem Grund ihres zuletzt belegten Literals der aktuellen Ebene, wodurch dieses Literal durch seine Ursachen ersetzt wird.
Sie endet, sobald genau ein Literal der aktuellen Ebene übrig ist.
Dieses Literal ist der \emph{first \acl{UIP}} (first \acs{UIP})\acused{UIP}: die dem Konflikt nächstgelegene Belegung, durch die jeder Pfad von der aktuellen Entscheidung zum Konflikt verläuft \vgl[280\psqq]{zhang2001}.
Die resultierende Klausel enthält den negierten \acs{UIP} und ansonsten nur Literale früherer Ebenen.
Da sie durch Resolution hergeleitet wird, folgt sie aus der Formel und ändert deren Erfüllbarkeit nicht.
Sie schließt die Kombination von Belegungen aus, die den Konflikt verursacht hat.

Anschließend führt der Solver einen \emph{Backjump} zur zweithöchsten Entscheidungsebene der gelernten Klausel aus, bzw. zu Ebene 0, wenn die Klausel nur ein Literal enthält.
Auf dieser Ebene sind alle ihre Literale außer dem negierten \acs{UIP} falsch, sodass die Klausel eine Unit-Klausel ist und sofort eine neue Belegung erzwingt.
Ebenen, deren Entscheidungen für den Konflikt keine Rolle spielten, werden vollständig übersprungen, anders als beim chronologischen Backtracking von DPLL.
Da jede gelernte Klausel eine Belegung ausschließt, die bereits zu einem Konflikt geführt hat, terminiert die Suche weiterhin.

Vier weitere Techniken entscheiden darüber, wie gut dies in der Praxis funktioniert.
Die Heuristik \emph{\acl{VSIDS}} (\acs{VSIDS})\acused{VSIDS} wählt die Entscheidungsvariablen: Variablen, die an jüngsten Konflikten beteiligt waren, erhalten einen höheren Score, und alle Scores klingen mit der Zeit ab, sodass der Solver dort verzweigt, wo die letzten Konflikte auftraten \vgl[531\psq]{moskewicz2001}.
\emph{Restarts} setzen die Suche in Abständen, die durch die Luby-Folge vorgegeben sind, auf Ebene 0 zurück, behalten aber die gelernten Klauseln und die Scores bei, damit eine ungünstige frühe Entscheidung den Solver nicht festhält \parencites[vgl.][532]{moskewicz2001}[178]{luby1993}.
Beim \emph{Phase Saving} erhält eine Entscheidungsvariable den Wert, den sie zuletzt hatte, statt eines festen Werts; dies führt die Suche nach einem Restart zu vielversprechenden Belegungen zurück \vgl[297]{pipatsrisawat2007}.
Schließlich kosten gelernte Klauseln Speicher und verlangsamen die Propagation, weshalb der Solver regelmäßig die weniger nützliche Hälfte von ihnen löscht.
Die Nützlichkeit wird über die \emph{\acl{LBD}} (\acs{LBD})\acused{LBD} gemessen, die Anzahl verschiedener Entscheidungsebenen in einer Klausel. Klauseln mit niedriger \acs{LBD} werden behalten \vgl[401\psq]{audemard2009}.

\subsubsection{DPLL vs.\ CDCL} \label{sec:comparison}

Tabelle~\ref{tab:dpll-vs-cdcl} fasst die Unterschiede zusammen.

\begin{table}[htb]
    \centering
    \caption{DPLL und \acs{CDCL} im Vergleich}
    \label{tab:dpll-vs-cdcl}
    \small
    \renewcommand{\arraystretch}{1.3}
    \begin{tabularx}{\textwidth}{@{}>{\raggedright\arraybackslash}p{3.6cm} >{\raggedright\arraybackslash}X >{\raggedright\arraybackslash}X@{}}
        \toprule
         & DPLL & \acs{CDCL} \\
        \midrule
        Reaktion auf einen Konflikt
            & versucht den entgegengesetzten Wert der letzten offenen Entscheidung
            & lernt eine Klausel, die durch Propagation eine neue Belegung erzwingt \\
        Backtracking
            & chronologisch, jeweils eine Ebene
            & nicht chronologisch, zur zweithöchsten Ebene der gelernten Klausel \\
        Gedächtnis für frühere Konflikte
            & keines
            & gelernte Klauseln, nach \acs{LBD} ausgedünnt \\
        Reihenfolge der Entscheidungen
            & Verzweigungsheuristik, unabhängig von Konflikten
            & \leavevmode\acs{VSIDS}, folgt jüngsten Konflikten \\
        Restarts
            & nein
            & ja, gelernte Klauseln bleiben erhalten \\
        \bottomrule
    \end{tabularx}
    \par\smallskip
    {\footnotesize Quelle: eigene Darstellung.}
\end{table}

Der Unterschied zeigt sich erst, wenn der Solver suchen muss.
Bei 20~$9 \times 9$-Rätseln mit eindeutiger Lösung und der minimalen Kodierung mussten beide Solver kaum suchen: DPLL benötigte im Median 24 Konflikte und 0,03\,s, \acs{CDCL} 9 Konflikte und 0,02\,s.
Von drei $16 \times 16$-Rätseln mit eindeutiger Lösung löste DPLL dagegen keines innerhalb des Zeitlimits von 20\,s, während \acs{CDCL} alle drei in einem Median von 1,3\,s löste.

\subsection{Implementierung} \label{sec:implementation}

Die beiden \acs{SAT}-Solver, DPLL und \acs{CDCL}, wurden ohne \acs{SAT}-Bibliothek in Python implementiert.
Abschnitt~\ref{sec:architecture} beschreibt die gemeinsame Architektur, die Abschnitte~\ref{sec:dpll-impl} und~\ref{sec:cdcl-impl} die Datenstrukturen der beiden Solver und Abschnitt~\ref{sec:trace} das Suchprotokoll, auf dem die Visualisierung aufbaut.

\subsubsection{Architektur} \label{sec:architecture}

Die Implementierung folgt den drei Schritten aus der Einleitung dieses Kapitels.
\texttt{encode} wandelt das Gitter in eine Liste von Klauseln und die Anzahl der Variablen um (Abschnitt~\ref{sec:encoding}).
Anschließend sucht einer der \acs{SAT}-Solver ein Modell, und \texttt{decode\_model} wandelt die wahren Variablen dieses Modells zurück in ein ausgefülltes Gitter.
Findet der Solver kein Modell, hat das Rätsel keine Lösung.

Alle \acs{SAT}-Solver teilen dieselbe Schnittstelle, \texttt{solve\_cnf(clauses, num\_vars)}.
Sie gibt ein Modell zurück, das für jede Variable genau ein Literal, $v$ oder $-v$, enthält, oder \texttt{None}, falls die Formel unerfüllbar ist, und lässt die Klauseln unverändert.
Zwei optionale Argumente sind allen Solvern gemeinsam: \texttt{stats} erfasst die Anzahl der Entscheidungen, Propagationen und Konflikte, und \texttt{trace} zeichnet jeden Suchschritt auf (Abschnitt~\ref{sec:trace}).
Das Aufzeichnen eines Protokolls verlangsamt einen Solver um 30\,\% bis hin zum Faktor 2,3.
Ohne Protokoll wird der Aufzeichnungscode übersprungen, und die verbleibenden Prüfungen kosten weniger als 5\,\% der Laufzeit, was innerhalb des Messrauschens liegt.
Dank dieser Schnittstelle lassen sich die Solver beliebig austauschen und an jeder \acs{KNF}-Formel testen, nicht nur an Sudoku.

Nichts im Code hängt von einer festen Gittergröße ab.
Jede Funktion leitet $n$ aus dem Gitter und die Blockgröße $b = \sqrt{n}$ aus $n$ ab, sodass derselbe Code $4 \times 4$-, $9 \times 9$- und $16 \times 16$-Gitter verarbeitet.

\subsubsection{Vom einfachen zum effizienten DPLL} \label{sec:dpll-impl}

Die erste DPLL-Implementierung folgt Abschnitt~\ref{sec:dpll} wörtlich.
Sie ist rekursiv, und jede Belegung erzeugt eine neue, vereinfachte Kopie der Formel ohne die erfüllten Klauseln und die falschen Literale.
Dadurch ist sie leicht zu überprüfen, doch jeder Suchschritt kostet Zeit proportional zur Größe der gesamten Formel.

Die effiziente Version führt dieselbe Suche mit anderen Datenstrukturen durch.
Statt die Formel zu kopieren, führt sie einen \emph{Trail}: die Liste der belegten Literale in der Reihenfolge, in der sie gesetzt wurden.
Backtracking kürzt den Trail einfach auf die Länge, die er vor der gescheiterten Entscheidung hatte.
Die offenen Entscheidungen liegen auf einem separaten Stack, sodass die Suche iterativ ist und nicht durch die Rekursionstiefe von Python begrenzt wird.

Die Unit Propagation verwendet \emph{Two Watched Literals} \vgl[531]{moskewicz2001}.
Jede Klausel beobachtet zwei ihrer Literale, die nicht falsch sind, und wird nur betrachtet, wenn eines davon falsch wird.
Der Solver sucht dann in der Klausel nach einem anderen nicht falschen Literal, das stattdessen beobachtet wird.
Gibt es keines, ist die Klausel zur Unit-Klausel oder zum Konflikt geworden.
Da Backtracking Belegungen nur aufhebt, bleiben die beobachteten Literale gültig und müssen nicht aktualisiert werden.

Entscheidungen folgen der Heuristik \texttt{shortest}: dem ersten freien Literal der kürzesten offenen Klausel, die nur positive Literale enthält.
Auf Sudoku ist dies die Zelle oder Einheit mit den wenigsten verbleibenden Kandidaten.
Die Beschränkung auf positive Klauseln ist wichtig: Die offenen binären \acs{AMO}-Klauseln sind stets die kürzesten, und eine Entscheidung \enquote{$x$ ist falsch} bewirkt wenig.

Sobald der Solver suchen muss, sparen diese Änderungen ein bis zwei Größenordnungen ein.
Auf dem leeren $16 \times 16$-Gitter mit der minimalen Kodierung benötigte die einfache Version 15,9\,s und die effiziente 0,36\,s.
\acs{CDCL} (Abschnitt~\ref{sec:cdcl-impl}) baut auf demselben Trail und denselben Watched Literals auf.

\subsubsection{CDCL} \label{sec:cdcl-impl}

Der \acs{CDCL}-Solver erweitert den effizienten DPLL.
Für jede Variable speichert er zusätzlich die Entscheidungsebene und den Grund, also die Klausel, die sie erzwungen hat. Entscheidungen haben keinen Grund.
Die Konfliktanalyse ist ein einziger Durchlauf rückwärts entlang des Trails (Listing~\ref{lst:analyze}).
Sie markiert die Variablen der Konfliktklausel und zählt, wie viele davon zur aktuellen Ebene gehören.
Anschließend nimmt sie das zuletzt belegte markierte Literal vom Trail und fährt, solange noch mehr als ein Literal der aktuellen Ebene offen ist, mit dem Grund dieses Literals fort.
Ist nur noch eines übrig, ist das gerade entnommene Literal der first \acs{UIP}.
Literale früherer Ebenen gehen direkt in die gelernte Klausel ein.
Die Funktion setzt den negierten \acs{UIP} an die erste und ein Literal der Backjump-Ebene an die zweite Stelle, damit diese beiden sofort beobachtet werden können, und berechnet im selben Durchlauf die \acs{LBD}.

% Pseudocode als Listing im Fließtext (nicht gleitend, Eintrag im Listingverzeichnis)
\par\medskip
\noindent\begin{minipage}{\linewidth}
\fcolorbox{black}{boxgray}{%
\begin{minipage}{\dimexpr\linewidth-2\fboxsep-2\fboxrule\relax}
\small
\begin{algorithmic}[1]
\Require Konfliktklausel $K_\bot$, deren Literale alle falsch sind
\Ensure gelernte Klausel $L$, Backjump-Ebene $b$, \acs{LBD} $\ell$
\Function{Analysiere}{$K_\bot$}
  \State $e \gets$ aktuelle Entscheidungsebene
  \State $L \gets [\,\cdot\,]$ \Comment{Platz 0 für das erzwingende Literal}
  \State $o \gets 0$ \Comment{offene Literale der Ebene $e$}
  \State $i \gets |T| - 1$ \Comment{$T$: Trail}
  \State $K \gets K_\bot$; \; $p \gets \text{nil}$
  \Repeat
    \ForAll{$q \in K$ mit $q \neq p$}
      \State $v \gets \operatorname{var}(q)$
      \If{$v$ nicht markiert \textbf{und} $\operatorname{ebene}(v) > 0$}
        \State markiere $v$; \; \Call{ErhöheScore}{$v$} \Comment{VSIDS}
        \If{$\operatorname{ebene}(v) = e$}
          \State $o \gets o + 1$
        \Else
          \State hänge $q$ an $L$ an
        \EndIf
      \EndIf
    \EndFor
    \While{$\operatorname{var}(T[i])$ nicht markiert} \Comment{zuletzt markiertes Literal}
      \State $i \gets i - 1$
    \EndWhile
    \State $p \gets T[i]$; \; $i \gets i - 1$
    \State entferne Markierung von $\operatorname{var}(p)$; \; $o \gets o - 1$
    \If{$o > 0$}
      \State $K \gets \operatorname{grund}(\operatorname{var}(p))$ \Comment{mit dem Grund von $p$ resolvieren}
    \EndIf
  \Until{$o = 0$} \Comment{$p$ ist der first \acs{UIP}}
  \State $L[0] \gets \neg p$
  \State entferne Markierungen aller Variablen in $L[1..]$
  \State $b \gets 0$
  \If{$|L| > 1$}
    \State $j \gets$ Index in $L[1..]$ mit höchster Entscheidungsebene \Comment{Backjump-Ebene}
    \State vertausche $L[1]$ und $L[j]$
    \State $b \gets \operatorname{ebene}(\operatorname{var}(L[1]))$
  \EndIf
  \State $\ell \gets |\{\operatorname{ebene}(\operatorname{var}(q)) : q \in L[1..]\}| + 1$
  \State \Return $(L, b, \ell)$
\EndFunction
\end{algorithmic}
\end{minipage}}
\captionof{lstlisting}{First-UIP-Konfliktanalyse als Pseudocode (nach \texttt{analyze} in \texttt{sat/cdcl.py})}
\label{lst:analyze}
{\footnotesize Quelle: eigene Darstellung.}
\end{minipage}
\par\medskip

Nach dem Backjump wird die gelernte Klausel zum Grund des neu erzwungenen Literals.
\acs{VSIDS} erhöht den Score jeder während der Analyse markierten Variable.
Statt nach jedem Konflikt alle Scores zu verringern, erhöht der Solver das Inkrement um den Faktor $1/0{,}95$, was dieselbe Wirkung hat.
Freie Variablen werden einem Heap entnommen, aus dem veraltete Einträge erst entfernt werden, wenn sie an die Reihe kommen.
Restarts erfolgen nach $100 \cdot \mathrm{luby}(i)$ Konflikten.
Ab 2.000 Konflikten löscht der Solver in wachsenden Abständen die Hälfte der gelernten Klauseln mit der höchsten \acs{LBD}, behält aber Klauseln mit einer \acs{LBD} von höchstens 2 sowie Klauseln, die gerade als Grund dienen.
Jede dieser Techniken lässt sich einzeln abschalten, sodass ihre Wirkung gemessen werden kann.

\subsubsection{Suchprotokoll} \label{sec:trace}

Um die Suche sichtbar zu machen, können beide Solver ein \emph{Protokoll} (\emph{Trace}) aufzeichnen: eine Liste mit einem Eintrag für jeden Suchschritt, die die Visualisierung anschließend Schritt für Schritt abspielen kann.
Eine Propagation wird mit dem Literal, seiner Entscheidungsebene und seinem Grund protokolliert.
Weitere Einträge kennzeichnen Entscheidungen, Konflikte und Backtracking, bei \acs{CDCL} zudem gelernte Klauseln, Restarts und das Löschen von Klauseln.

Auf Sudoku übertragen entspricht jeder Eintrag direkt einer Änderung des Gitters.
Eine auf wahr gesetzte Variable $x_{r,c,d}$ trägt die Ziffer $d$ in Zelle $(r, c)$ ein, eine auf falsch gesetzte entfernt $d$ aus den Kandidaten dieser Zelle.
Die meisten Schritte sind von der zweiten Art: In einem \acs{CDCL}-Protokoll des schweren $9 \times 9$-Rätsels \enquote{AI Escargot} entfernen 82\,\% aller Einträge einen Kandidaten.
Ein \acs{SAT}-Solver lässt sich daher am besten auf einem Gitter mit Kandidaten darstellen, nicht als Folge gesetzter Ziffern.
Mit der minimalen Kodierung umfasst das Protokoll von AI Escargot 1.383 Einträge für DPLL und 11.449 für \acs{CDCL}.
